%% kapa: exam=1 round=25 year=2014 period=1 form=A subject=civil_law subject_ko=민법 seq=29 number=29
%% source: Q-Net 감정평가사 기출문제 (cst00309) artlSeq=5101207 file=extracted/1008/5101207/2014년 제25회 감정평가사 1차시험 기출문제/1교시(A형).hwp
%% conversion: kordoc 4.24.51 -> tools/gen1.py (auto); review=proofread-sonnet; flags=none
\begin{question}{29}
상린관계(相隣關係)에 관한 설명으로 옳지 않은 것은?
\begin{choices}
  \item 건물을 축조함에는 특별한 관습이 없으면, 경계로부터 반 미터 이상의 거리를 두어야 한다.
  \item 토지소유자는 이웃 토지로부터 자연히 흘러오는 물을 막지 못한다.
  \item 인접하여 토지를 소유한 자는 다른 관습이 없으면, 쌍방이 공동비용으로 통상의 경계표를 설치하기 위한 측량비용을 절반하여 부담한다.
  \item 고지(高地)의 소유자는 이웃 저지(低地)에 자연히 흘러내리는 이웃 저지에서 필요한 물을 자기의 정당한 사용범위를 넘어서 이를 막지 못한다.
  \item 흐르는 물이 저지에서 폐색(閉塞)된 때에는 특별한 관습이 없으면, 고지소유자는 자비로 소통에 필요한 공사를 할 수 있다.
\end{choices}
\end{question}
\answer{3}
